\section{Complete proof of the mixed-parity \texorpdfstring{$60$}{60} theorem}\label{app:mixed60}

Let $a,b,c$ be odd and $e$ even.  We prove that these four leaves can be
used exactly once to construct a multiple of $60$.  The top-level partition
is exhaustive: Case A handles an input divisible by $5$; after excluding
that case, Case B handles $3\mid e$; under $3\nmid e$, Cases D, E, and C
handle respectively zero, exactly one, and at least two odd multiples of
$3$.  Within every case, a statement that one expression is divisible by two
moduli always refers to one and the same choice of sign.  Signs belonging to
different disjoint blocks may be chosen independently; within any single
block, however, the same sign is used simultaneously for all stated
congruence conditions.

\paragraph{Case A: one input is divisible by $5$.}
If an odd input, say $a$, is divisible by $5$, multiply it by a
$12$-multiple constructed from $b,c,e$.  If $5\mid e$, then $10\mid e$ and
the remaining three odd inputs yield a $12$-multiple.  Hence this case is
complete.  From now on all four inputs are nonzero modulo $5$.

\paragraph{Case B: $3\mid e$.}
Two of the three odd inputs lie in one signed class modulo $5$, so their
suitable sum or difference is divisible by $5$ and, because both are odd, by
$2$; this is a $10$-block.  The remaining odd input multiplied by $e$ is
divisible by $6$.  Hence $10\cdot6=60$.  From now on $3\nmid e$.

\paragraph{Case C: at least two odd inputs are divisible by $3$.}
If two such odd inputs lie in one signed class modulo $5$, their suitable sum
or difference is divisible by $2$, $3$, and $5$, hence by $30$; the
remaining odd-even product is even.

Otherwise exactly two odd inputs are divisible by $3$ and they occupy the
two different signed classes modulo $5$.  The third odd input lies in the
same class as one of them.  That same-class odd pair gives a $10$-block,
while the unused multiple-of-$3$ odd input multiplied by $e$ gives a
$6$-block.  Thus this case is complete.

\paragraph{Case D: no odd input is divisible by $3$.}
Now all four inputs are coprime to $15$.  Choose two odd inputs $a,b$ in one
signed class modulo $5$.  We distinguish the class positions of $c,e$.

\smallskip
\noindent\emph{D1. The remaining inputs $c,e$ occupy one signed class modulo
$5$.}
For the pair $a,b$, compare the three signs
$\sigma_3(a,b),\sigma_4(a,b),\sigma_5(a,b)$.  At least two are equal.
If $\sigma_3=\sigma_5$, the pair yields $30$ and $ce$ yields $2$.  If
$\sigma_4=\sigma_5$, the pair yields $20$, while $c\pm e$ can be made
divisible by $3$.  If $\sigma_3=\sigma_4$, the pair yields $12$, while the
same-class pair $c,e$ yields $5$.  Each branch gives $60$.

\smallskip
\noindent\emph{D2. The three odd inputs $a,b,c$ occupy one signed class
modulo $5$, and $e$ occupies the other.}
If an odd pair has $\sigma_3=\sigma_5$, then the same signed sum or
difference is divisible by $2$, $3$, and $5$, hence forms a $30$-block; the
remaining odd-even product supplies $2$.  If an odd pair has
$\sigma_4=\sigma_5$, then the same signed sum or difference forms a
$20$-block, while the remaining odd-even pair can be combined into a
$3$-block.

Suppose neither exit occurs for any odd pair.  Since each admissible sign is
binary, every edge of the odd triangle must satisfy
\[
 \sigma_3=\sigma_4=-\sigma_5.
\]
The signed-class triangle lemma says that the product of the three
$\sigma_3$ edge signs is $-1$ and that the product of the three $\sigma_5$
edge signs is also $-1$.  On the other hand, multiplying the three relations
$\sigma_5=-\sigma_3$ shows that the $\sigma_5$ edge-sign product is the
negative of the $\sigma_3$ edge-sign product, because three negative factors
occur.  The $\sigma_3$ product is $-1$, so this would make the $\sigma_5$
product $+1$, a contradiction.  Hence one of the two preceding exits must
occur.

\smallskip
\noindent\emph{D3. The inputs $a,b,e$ occupy one signed class modulo $5$,
and $c$ the other.}
If $\sigma_3(a,b)=\sigma_5(a,b)$, then the same signed sum or
difference of $a,b$ forms a $30$-block and the complementary product $ce$
supplies $2$.  If $\sigma_4(a,b)=\sigma_5(a,b)$, then the same signed sum
or difference of $a,b$ forms a $20$-block and the complementary pair $c,e$
can be combined into a $3$-block.  Otherwise put
\[
 \rho(u,v)=\sigma_3(u,v)\sigma_5(u,v).
\]
On the same-class triangle $a,b,e$, the two triangle identities give
\[
 \rho(a,b)\rho(b,e)\rho(e,a)=1.
\]
Since $\rho(a,b)=-1$, exactly one of the other two factors equals $+1$.
Thus either $a,e$ or $b,e$ has one sum or difference divisible by both $3$
and $5$, giving a $15$-block.  The complementary odd pair yields $4$.

This completes the case in which no odd input is divisible by $3$.

\paragraph{Case E: exactly one odd input is divisible by $3$.}
Let $3\mid a$ and $3\nmid bce$.  If $b,c$ occupy one signed class modulo
$5$, they yield a $10$-block and $ae$ yields a $6$-block.  Hence assume
$b,c$ occupy different signed classes.  Relabel so that $a,b$ occupy one
class and $c$ the other.  There are two positions for $e$.

\subsection*{\texorpdfstring{E1. The even input is in the class of $a,b$}{E1. The even input is in the class of a,b}}\label{app:caseE1}

By sign normalization, choose representatives satisfying
\[
 a\equiv b\equiv e\equiv r,
 \qquad c\equiv-2r\pmod5.
\]
Consequently
\[
\begin{gathered}
5\mid a-b,\quad 5\mid b-e,\quad5\mid a-e,\\
5\mid a+b+c,\quad5\mid a+c+e,\quad5\mid b+c+e.
\end{gathered}
\]
If $4\mid a-b$, then $a-b$ is a $20$-block and $c,e$ can be combined
into a $3$-block.  If $3\mid b-e$, then $b-e$ is a $15$-block and $a,c$
can be combined into a $4$-block.  If there exists a single sign
$s\in\{\pm1\}$ such that
\[
 3\mid b+sc
 \qquad\text{and}\qquad
 4\mid b+sc,
\]
then $b+sc$ is a $12$-block and $a-e$ is a $5$-block.

Assume none of these three exits occurs.  Since $a,b$ are odd, exactly one
of $a-b$ and $a+b$ is divisible by $4$, so $4\nmid a-b$ implies
$4\mid a+b$.  Since $b,e$ are both nonzero modulo $3$, exactly one of
$b-e$ and $b+e$ is divisible by $3$, so $3\nmid b-e$ implies
$3\mid b+e$.  Finally, for the odd pair $b,c$, the modulo-$3$ sign and the
modulo-$4$ sign are each unique.  The absence of a single sign producing a
$12$-block means that these two signs are opposite.  Thus
\[
 4\mid a+b,
 \qquad 3\mid b+e,
\]
and the signs that make $b\pm c$ divisible by $3$ and by $4$ are opposite.

First let $e\equiv0\pmod4$ and put $T=b+c+e$.  We know $5\mid T$.  If
$4\mid T$, then $aT$ is divisible by $3\cdot20=60$.  If $4\nmid T$, then
$b+c\equiv2\pmod4$, so $4\mid b-c$.  Because the modulo-$3$ and modulo-$4$
signs for $b\pm c$ are opposite, $3\mid b+c$.  Hence
$15\mid a+b+c$, and multiplication by the unused input $e$, divisible by
$4$, gives $60$.

Now let $e\equiv2\pmod4$, again with $T=b+c+e$.  If $4\mid T$, the product
$aT$ is divisible by $60$.  Otherwise $4\mid b+c$, so the opposite-sign
condition gives $3\mid b-c$.  Together with $3\mid b+e$, this yields
$3\mid c+e$.  Put
\[
 U=a+c+e.
\]
Modulo $5$, $U\equiv r-2r+r\equiv0$.  Modulo $3$, both $a$ and $c+e$ are
zero.  Finally $4\mid a+b$ and $4\mid b+c$ imply
$a\equiv c\equiv-b\pmod4$, and therefore
\[
 U\equiv-2b+2\equiv0\pmod4
\]
because $b$ is odd.  Thus $60\mid U$, and $bU$ uses the final leaf.

\subsection*{\texorpdfstring{E2. The even input is in the class of $c$}{E2. The even input is in the class of c}}\label{app:caseE2}

Normalize signs so that
\[
 a\equiv b\equiv r,
 \qquad c\equiv e\equiv-2r\pmod5.
\]
Then
\[
5\mid a-b,\quad5\mid c-e,\quad
5\mid b-c-e,\quad5\mid a-c-e.
\]
If $4\mid a-b$, combine the resulting $20$-block with a $3$-block from
$c,e$.  If $3\mid c-e$, combine the resulting $15$-block with a $4$-block
from $a,b$.  Assume neither exit occurs.  Since $a,b$ are odd, exactly one
of $a-b$ and $a+b$ is divisible by $4$; hence $4\nmid a-b$ implies
$4\mid a+b$.  Since $c,e$ are nonzero modulo $3$, exactly one of $c-e$ and
$c+e$ is divisible by $3$; hence $3\nmid c-e$ implies $3\mid c+e$.  Thus
\[
4\mid a+b,
\qquad3\mid c+e.
\]
Put $T=b-c-e$.  If $4\mid T$, then $aT$ is divisible by $60$.  Otherwise
$T$ is even but not divisible by $4$, so $T\equiv2\pmod4$.  Put
$U=a-c-e$.  We have $5\mid U$ from the normalization and $3\mid U$ from
$3\mid a$ and $3\mid c+e$.  Moreover $a\equiv-b\pmod4$ and
$c+e\equiv b-2\pmod4$, so
\[
 U\equiv-b-(b-2)=2-2b\equiv0\pmod4
\]
because $b$ is odd.  Thus $60\mid U$, and $bU$ uses the remaining input.

Cases A--E exhaust all possibilities.  Therefore
$3\OO1\EE\to60$.
